\documentclass[11pt]{article}
\usepackage{fontspec}
\setmainfont[
  Path=fonts/,
  Extension=.ttf,
  UprightFont=*-400,
  BoldFont=*-700
]{NotoSerifSC}
\setsansfont[
  Path=fonts/,
  Extension=.ttf,
  UprightFont=*-400,
  BoldFont=*-700
]{NotoSerifSC}
\XeTeXlinebreaklocale "zh"
\XeTeXlinebreakskip=0pt plus 1pt
\usepackage[a4paper,margin=1.70cm]{geometry}
\usepackage{amsmath,amssymb,mathtools,bm}
\usepackage{booktabs,longtable,array}
\usepackage{enumitem}
\setlength{\parindent}{2em}
\setlength{\parskip}{0.35em}
\setlist{nosep,leftmargin=2em}
\allowdisplaybreaks[3]
\numberwithin{equation}{section}
\pagestyle{plain}

\newcommand{\R}{\mathbb{R}}
\newcommand{\E}{\mathbb{E}}
\newcommand{\Pp}{\mathbb{P}}
\newcommand{\1}{\mathbf{1}}
\newcommand{\T}{\mathsf{T}}
\newcommand{\F}{\mathrm{F}}
\newcommand{\cN}{\mathcal{N}}
\newcommand{\cL}{\mathcal{L}}
\newcommand{\cS}{\mathcal{S}}
\newcommand{\cH}{\mathcal{H}}
\newcommand{\cA}{\mathcal{A}}
\newcommand{\cM}{\mathcal{M}}
\newcommand{\cT}{\mathcal{T}}
\newcommand{\diag}{\operatorname{diag}}
\newcommand{\tr}{\operatorname{tr}}
\newcommand{\Var}{\operatorname{Var}}
\newcommand{\Cov}{\operatorname{Cov}}
\newcommand{\softmax}{\operatorname{softmax}}
\newcommand{\KL}{D_{\mathrm{KL}}}
\newcommand{\sg}{\operatorname{sg}}
\begin{document}
    \begin{center}
      {\LARGE\bfseries 4-2103.00020v1-Visual-CLIP}
    \end{center}
    \vspace{0.8em}

    \section{符号与维度}
    \begin{longtable}{>{$}l<{$} >{$}l<{$} p{10cm}}
    \toprule
    \text{符号} & \text{维度} & \text{含义}\\
    \midrule
    u_i,v_i & \R^d & 第 $i$ 个图像与文本的单位化嵌入。\\
        s_{ij} & \R & 图像 $i$ 与文本 $j$ 的缩放余弦 logit。\\
        \tau & \R_{>0} & softmax 温度。\\
        N & \mathbb N & 对比学习批大小。\\
    \bottomrule
    \end{longtable}

    所有向量默认是列向量；未特别说明时，期望同时对数据、噪声和采样时刻取值。下文只展开论文中承担方法逻辑的公式，并把所需假设写在推导发生的位置。

\section{对称图文对比目标}

一批 $N$ 对图像与文本，经编码和单位化得到
$u_i,v_j\in\R^d$，$\|u_i\|=\|v_j\|=1$。logit
\begin{equation}
 z_{ij}=\frac{u_i^Tv_j}{\tau}.
\end{equation}
图到文损失
\begin{equation}
 \cL_{I\to T}=-\frac1N\sum_i
 \log\frac{e^{z_{ii}}}{\sum_je^{z_{ij}}},
\end{equation}
文到图损失对列作 softmax：
\begin{equation}
 \cL_{T\to I}=-\frac1N\sum_j
 \log\frac{e^{z_{jj}}}{\sum_ie^{z_{ij}}}.
\end{equation}
总损失常取二者平均。对一行，记
$p_{ij}=e^{z_{ij}}/\sum_ke^{z_{ik}}$，则
\begin{equation}
 \frac{\partial\cL_{I\to T}}{\partial z_{ij}}
 =\frac1N(p_{ij}-\1\{i=j\}).
\end{equation}
因此
\begin{equation}
 \frac{\partial\cL_{I\to T}}{\partial u_i}
 =\frac1{N\tau}\left(\sum_jp_{ij}v_j-v_i\right).
\end{equation}
梯度把图像向量拉向正文本，并按当前混淆概率推离负文本；困难负样本权重更大。

\section{单位化的 Jacobian}

网络原始输出 $h\ne0$，单位向量 $u=h/\|h\|_2$。令 $r=\|h\|$，
\begin{equation}
 du=\frac{dh}{r}-\frac{h(h^Tdh)}{r^3}
 =\frac1r(I-uu^T)dh.
\end{equation}
故
\begin{equation}
 \boxed{J_{u\leftarrow h}=\frac1{\|h\|}(I-uu^T).}
\end{equation}
反向梯度
\begin{equation}
 \nabla_h\cL=\frac1{\|h\|}(I-uu^T)\nabla_u\cL
\end{equation}
与 $u$ 正交，所以一阶更新只改变方向而不改变径向分量。若 $\|h\|$ 很小，
梯度会放大，实际实现需用数值下界。

单位化后相似度是余弦，$-1\le u^Tv\le1$，模型不能靠增大特征范数无限降低交叉熵；
整体 logit 尺度交给温度参数。

\section{温度的导数与熵效应}

令 $s_{ij}=u_i^Tv_j$、$z_{ij}=s_{ij}/\tau$。单行损失
\begin{equation}
 \ell_i=-\frac{s_{ii}}\tau+\log\sum_j e^{s_{ij}/\tau}.
\end{equation}
对 $\tau$ 求导：
\begin{equation}
 \frac{\partial\ell_i}{\partial\tau}
 =\frac1{\tau^2}\left(s_{ii}-\sum_jp_{ij}s_{ij}\right).
\end{equation}
若正相似度高于当前 softmax 平均，梯度下降会减小 $\tau$，使分布更尖锐。
实践中常参数化 $\gamma=\log(1/\tau)$，此时
\begin{equation}
 \frac{\partial\ell_i}{\partial\gamma}
 =\sum_jp_{ij}z_{ij}-z_{ii}.
\end{equation}
对 $\gamma$ 限幅可防止温度趋近零造成 logit 与梯度过大。

\section{InfoNCE 互信息下界的条件}

设正样本 $(X,Y_1)\sim p(x,y)$，其余 $Y_2,\ldots,Y_N$ 独立来自边缘 $p(y)$。
从 $N$ 个候选中识别正样本索引 $I$。最优分类后验正比密度比
\begin{equation}
 \Pp(I=i\mid x,y_{1:N})
 \propto\frac{p(y_i\mid x)}{p(y_i)}.
\end{equation}
用得分 $f(x,y)>0$ 近似该密度比，InfoNCE
\begin{equation}
 \cL_N=-\E\log\frac{f(X,Y_1)}{\sum_{j=1}^{N}f(X,Y_j)}.
\end{equation}
标准推导给出
\begin{equation}
 \boxed{I(X;Y)\ge\log N-\cL_N.}
\end{equation}
它要求负样本来自正确边缘分布并独立；批内语义重复、配对噪声或数据相关会改变该严格解释。
此外下界至多为 $\log N$，大批量提高可表达的下界上限，但计算和假负样本概率也上升。

\section{批内负样本梯度的统计耦合}

每个样本既是自身正样本，又是其他 $N-1$ 个样本的负样本。因此 $v_j$ 的总梯度含
\begin{equation}
 \nabla_{v_j}\cL_{I\to T}
 =\frac1{N\tau}\left(\sum_i p_{ij}u_i-u_j\right).
\end{equation}
若两个不同索引实际语义相同，项 $p_{ij}u_i$ 仍把它当负样本推开，这就是假负样本偏差。
增大批量降低 softmax 分母的抽样噪声，却可能增加语义碰撞，二者需同时考虑。

若使用跨设备 all-gather，只有本设备特征保留梯度而远端特征停止梯度，则前向负样本集合相同，
但远端样本不接收作为 negative 的反向贡献；这与全局可微 gather 的优化目标并不完全相同。

\section{动量编码器与软目标}

动量参数递推
\begin{equation}
 \bar\theta_k=m\bar\theta_{k-1}+(1-m)\theta_k
\end{equation}
展开为
\begin{equation}
 \bar\theta_k=m^k\bar\theta_0+(1-m)
 \sum_{j=1}^{k}m^{k-j}\theta_j.
\end{equation}
因此教师是历史学生的指数平均，有效窗口约为 $1/(1-m)$。

若硬标签 $e_i$ 与动量教师分布 $q_i$ 混合
\begin{equation}
 y_i=(1-\alpha)e_i+\alpha q_i,
\end{equation}
交叉熵对学生 logit 的梯度为
\begin{equation}
 \nabla_z\ell=p_i-y_i
 =(1-\alpha)(p_i-e_i)+\alpha(p_i-q_i).
\end{equation}
所以软蒸馏项把教师对其他样本的相似性结构传给学生，同时可能传播教师偏差。

\section{图文匹配与对比学习的差别}

双塔对比模型的得分是独立编码后的内积，所有图文表征可预计算；跨模态融合模型让 token 之间交互，
匹配 logit $a_\theta(I,T)$ 的二元损失为
\begin{equation}
 \ell_{\rm ITM}=-y\log\sigma(a)-(1-y)\log(1-\sigma(a)),
\qquad
 \frac{\partial\ell}{\partial a}=\sigma(a)-y.
\end{equation}
它能表达细粒度对应，但每一候选对都需联合前向。实践中先用对比相似度检索困难负样本，
再用融合编码器判别，是计算量与交互能力之间的分解。

\section{自回归似然与 token 梯度}

对序列 $x_{1:T}$，因果语言模型按概率链式法则分解
\begin{equation}
 p_\theta(x_{1:T})=\prod_{t=1}^{T}p_\theta(x_t\mid x_{<t}).
\end{equation}
负对数似然
\begin{equation}
 \cL_{\rm NLL}=-\sum_{t=1}^{T}\log p_\theta(x_t\mid x_{<t}).
\end{equation}
若第 $t$ 步 logits 为 $z_t\in\R^{V}$，概率
$p_{tj}=e^{z_{tj}}/\sum_ke^{z_{tk}}$，真实 token 索引 $y_t$，则
\begin{equation}
 \ell_t=-z_{t,y_t}+\log\sum_ke^{z_{tk}},
\end{equation}
逐 logit 梯度
\begin{equation}
 \frac{\partial\ell_t}{\partial z_{tj}}
 =p_{tj}-\1\{j=y_t\}.
\end{equation}
Hessian 正是 softmax Jacobian
\begin{equation}
 \nabla_z^2\ell_t=\diag(p_t)-p_tp_t^T\succeq0.
\end{equation}
所以交叉熵关于 logits 凸，但关于深网参数并不凸。

对非 padding token 集合 $\mathcal I$，平均损失和困惑度
\begin{equation}
 \bar L=\frac1{|\mathcal I|}\sum_{t\in\mathcal I}\ell_t,
 \qquad \operatorname{PPL}=e^{\bar L}.
\end{equation}
因此 PPL 是真实 token 概率倒数的几何平均：
\begin{equation}
 \operatorname{PPL}=\left(
 \prod_{t\in\mathcal I}\frac1{p_\theta(x_t\mid x_{<t})}
 \right)^{1/|\mathcal I|}.
\end{equation}
不同 tokenizer 的 token 粒度不同，PPL 不能不经换算直接横向比较。

\section{旋转位置编码的相对位置恒等式}

对二维坐标对，定义旋转
\begin{equation}
 R(m\omega)=
 \begin{bmatrix}
 \cos(m\omega)&-\sin(m\omega)\\
 \sin(m\omega)&\cos(m\omega)
 \end{bmatrix}.
\end{equation}
位置 $m,n$ 的 query、key 分别为
$q_m=R(m\omega)q$、$k_n=R(n\omega)k$。内积
\begin{align}
 q_m^Tk_n
 &=q^TR(m\omega)^TR(n\omega)k\\
 &=q^TR((n-m)\omega)k.
\end{align}
第二步用到 $R(a)^T=R(-a)$ 和 $R(a)R(b)=R(a+b)$。
所以绝对位置旋转后的内积只依赖相对距离 $n-m$。

在 $d_h$ 维中对相邻坐标对使用频率
\begin{equation}
 \omega_j=\theta^{-2j/d_h},\qquad j=0,\ldots,d_h/2-1.
\end{equation}
低 $j$ 频率高、分辨近距离，高 $j$ 频率低、覆盖远距离。若推理长度远超训练长度，
相位 $m\omega_j$ 的外推和周期混叠会改变注意力几何；位置插值本质上是重新缩放这些相位。

\section{RMSNorm 的 Jacobian}

忽略可学习增益，RMSNorm
\begin{equation}
 y=\frac{x}{s},\qquad
 s=\sqrt{\frac1d x^Tx+\epsilon}.
\end{equation}
微分
\begin{equation}
 ds=\frac{x^Tdx}{d s},
\end{equation}
因此
\begin{align}
 dy&=\frac1s dx-\frac{x}{s^2}ds\\
 &=\left(\frac1sI-\frac{xx^T}{d s^3}\right)dx.
\end{align}
故
\begin{equation}
 J_{\rm RMS}=\frac1sI-\frac{xx^T}{d s^3}.
\end{equation}
与 LayerNorm 相比，它没有去均值投影
$I-d^{-1}\1\1^T$，因此不消除常数方向。加上逐坐标增益 $g$ 后，
Jacobain 左乘 $\diag(g)$。

谱上，对任意 $v\perp x$，$Jv=v/s$；沿 $x$ 方向
\begin{equation}
 Jx=\left(\frac1s-\frac{\|x\|^2}{d s^3}\right)x
 =\frac{\epsilon}{s^3}x.
\end{equation}
当 $\epsilon$ 很小，径向尺度方向被强烈压制，而切向方向按 $1/s$ 缩放。

\section{SwiGLU 前馈层的逐项梯度}

SwiGLU 可写
\begin{equation}
 a=xW_a,\qquad b=xW_b,\qquad
 h=\operatorname{swish}(a)\odot b,\qquad
 y=hW_o,
\end{equation}
其中
\begin{equation}
 \operatorname{swish}(a)=a\sigma(a).
\end{equation}
其导数
\begin{align}
 \frac d{da}[a\sigma(a)]
 &=\sigma(a)+a\sigma(a)(1-\sigma(a)).
\end{align}
令 $g_h=\partial\cL/\partial h$，则
\begin{align}
 g_a&=g_h\odot b\odot
 [\sigma(a)+a\sigma(a)(1-\sigma(a))],\\
 g_b&=g_h\odot\operatorname{swish}(a),\\
 \nabla_{W_a}\cL&=x^Tg_a,\qquad
 \nabla_{W_b}\cL=x^Tg_b,\\
 \nabla_x\cL&=g_aW_a^T+g_bW_b^T.
\end{align}
门分支 $b$ 同时调制激活和 $a$ 分支梯度，故两个投影不是可独立解释的普通 MLP。

\section{MoE 路由与负载}

对 token $x_i$，router logits $r_i=W_rx_i$，概率
\begin{equation}
 p_{ie}=\frac{e^{r_{ie}}}{\sum_{j=1}^{E}e^{r_{ij}}}.
\end{equation}
若 top-$k$ 专家集合 $S_i$，归一化门
\begin{equation}
 \widetilde p_{ie}=\frac{p_{ie}\1\{e\in S_i\}}
 {\sum_{j\in S_i}p_{ij}},
 \qquad
 y_i=\sum_{e\in S_i}\widetilde p_{ie}F_e(x_i).
\end{equation}
固定集合 $S_i$ 内可正常求导；集合边界的 arg-top-$k$ 不连续，几乎处处把未选专家梯度置零。

令软重要性
\begin{equation}
 P_e=\frac1N\sum_ip_{ie},
\end{equation}
硬负载
\begin{equation}
 F_e=\frac1N\sum_i\1\{e\in S_i\}/k.
\end{equation}
常见平衡项
\begin{equation}
 L_{\rm bal}=E\sum_eP_eF_e
\end{equation}
在均匀 $P_e=F_e=1/E$ 时等于 $1$。把硬 $F_e$ 停止梯度后，
router 梯度只通过 $P_e$，会降低当前过载专家的概率；它是训练启发式，不等同于严格容量约束。

若每专家容量为
\begin{equation}
 C_e=\left\lceil\frac{Nk}{E}\,c_{\rm cap}\right\rceil,
\end{equation}
超过容量的 token 必须丢弃或转送。即使平均负载均匀，局部批次方差仍可能溢出。

\section{低秩 KV 潜变量的吸收}

设 key/value 由共享潜变量
\begin{equation}
 c_t=W_cx_t\in\R^{d_c},\qquad
 k_t=W_kc_t,\qquad v_t=W_vc_t.
\end{equation}
注意力 logit
\begin{equation}
 q_s^Tk_t=q_s^TW_kc_t=(W_k^Tq_s)^Tc_t.
\end{equation}
因此解码时可把 $W_k$ 吸收到 query 投影，缓存 $c_t$ 而非完整 $k_t$。
若 value 聚合
\begin{equation}
 o_s=\sum_tp_{st}W_vc_t=W_v\left(\sum_tp_{st}c_t\right),
\end{equation}
$W_v$ 可放到加权和之后。缓存从每 token 的 $d_k+d_v$ 降为 $d_c$，
前提是位置编码等操作不破坏这些线性交换关系；对只作用在部分 key 维度的 RoPE，需把旋转部分另行缓存或重构。

\section{多 token 预测的联合似然}

若同一隐藏状态 $h_t$ 用 $K$ 个头预测 $x_{t+1},\ldots,x_{t+K}$，
\begin{equation}
 \cL_t=\sum_{k=1}^{K}\lambda_k
 [-\log p_{\theta,k}(x_{t+k}\mid h_t)].
\end{equation}
共享骨干梯度
\begin{equation}
 \nabla_{h_t}\cL_t=\sum_{k=1}^{K}\lambda_k g_{t,k}.
\end{equation}
其平方范数精确展开
\begin{equation}
 \left\|\sum_k\lambda_kg_{t,k}\right\|^2
 =\sum_k\lambda_k^2\|g_{t,k}\|^2
 +2\sum_{j<k}\lambda_j\lambda_k g_{t,j}^Tg_{t,k}.
\end{equation}
交叉内积为负时存在任务冲突；辅助头退火相当于训练后期逐渐去掉这些交叉项，
使最终目标回到下一 token 似然。

\end{document}
